% === ISTILAH ===

\newglossaryentry{finetuning}{
    name={\textit{fine-tuning}},
    description={Proses penyesuaian lanjutan model terlatih (\textit{pre-trained}) pada data tugas atau domain tertentu agar berkinerja lebih baik pada tugas tersebut}
}

\newglossaryentry{reranking}{
    name={\textit{re-ranking}},
    description={Tahap penyusunan ulang peringkat dokumen/pasase hasil pencarian tahap awal agar dokumen yang paling relevan menempati posisi teratas sebelum diberikan ke model generatif}
}

\newglossaryentry{workflow}{
    name={\textit{workflow automation}},
    description={Otomatisasi rangkaian tugas atau proses kerja yang sebelumnya dilakukan manual, pada konteks ini dijalankan oleh agen berbasis model bahasa}
}

% === SINGKATAN ===

\newacronym{ai}{AI}{\textit{Artificial Intelligence}}
\newacronym{nlp}{NLP}{\textit{Natural Language Processing}}
\newacronym{llm}{LLM}{\textit{Large Language Model}}
\newacronym{slm}{SLM}{\textit{Small Language Model}}
\newacronym{lora}{LoRA}{\textit{Low-Rank Adaptation}}
\newacronym{peft}{PEFT}{\textit{Parameter-Efficient Fine-Tuning}}
\newacronym{rag}{RAG}{\textit{Retrieval-Augmented Generation}}
\newacronym{api}{API}{\textit{Application Programming Interface}}
\newacronym{gpu}{GPU}{\textit{Graphics Processing Unit}}
